\documentclass[11pt,a4paper]{article}
\usepackage[margin=18mm]{geometry}
\usepackage[T1]{fontenc}
\usepackage{booktabs}
\setlength{\parindent}{0pt}
\setlength{\parskip}{7pt}
\pagestyle{plain}

\begin{document}

{\Large\bfseries Practical 9: Advanced Sorting}\par
Student ID: 23472297

\textbf{Runtime results}

The table contains measured runtimes from \texttt{python3 benchmark.py}.
Each case runs three times: the first is discarded and the remaining two
are averaged. Only sorting is timed; array generation, printing and
validation are excluded. Fixed seeds give each algorithm the same input
arrays for each repeat. All times are in milliseconds and may vary between runs.

\begin{center}
\small
\setlength{\tabcolsep}{4pt}
\begin{tabular}{@{}rlrrrrrrr@{}}
\toprule
Size & Input & Bubble & Insertion & Selection & Merge & Q-left & Q-med & Q-rand \\
\midrule
250 & Asc. & 0.019 & 0.019 & 2.195 & 0.371 & 1.242 & 0.164 & 0.226 \\
250 & Desc. & 5.451 & 5.013 & 2.284 & 0.350 & 2.493 & 0.272 & 0.236 \\
250 & Random & 3.658 & 2.540 & 2.933 & 0.402 & 0.226 & 0.215 & 0.222 \\
\midrule
500 & Asc. & 0.048 & 0.045 & 9.334 & 0.787 & 4.881 & 0.405 & 0.538 \\
500 & Desc. & 22.200 & 20.401 & 8.792 & 0.802 & 10.887 & 0.700 & 0.593 \\
500 & Random & 16.373 & 10.110 & 8.669 & 0.925 & 0.505 & 0.501 & 0.658 \\
\midrule
1000 & Asc. & 0.101 & 0.078 & 40.741 & 1.777 & 20.183 & 1.376 & 1.245 \\
1000 & Desc. & 110.538 & 89.967 & 36.998 & 1.918 & 44.832 & 1.693 & 1.259 \\
1000 & Random & 67.844 & 43.046 & 37.931 & 2.178 & 1.308 & 1.257 & 1.324 \\
\midrule
2000 & Asc. & 0.165 & 0.153 & 155.343 & 4.027 & 77.748 & 2.081 & 2.801 \\
2000 & Desc. & 392.353 & 343.776 & 143.513 & 3.976 & 185.991 & 3.945 & 2.989 \\
2000 & Random & 276.884 & 183.436 & 139.538 & 4.596 & 2.621 & 2.521 & 2.786 \\
\bottomrule
\end{tabular}
\end{center}

Q-left, Q-med and Q-rand mean QuickSort with leftmost, median-of-three
and random pivots, respectively. Asc. means ascending; Desc. means descending.

\textbf{Discussion}

Bubble and Insertion Sort were fastest on ascending input, matching their
$O(n)$ best case. On random and descending input, their runtimes increased
much faster, consistent with $O(n^2)$ behaviour. Selection Sort has
$O(n^2)$ time in all cases and did not benefit much from sorted input.
Merge Sort performed consistently across input conditions, with
$O(n\log n)$ time in every case, but requires $O(n)$ extra array space.
It is stable because equal values are taken from the left half first.
Leftmost-pivot QuickSort was much slower on ascending and descending input
because repeatedly unbalanced partitions produce $O(n^2)$ time.
Median-of-three and random pivots improved these cases substantially.
All QuickSort versions were faster than Merge Sort on the tested random
inputs. QuickSort has average or expected $O(n\log n)$ time, but all three
versions can still take $O(n^2)$ in the worst case. It rearranges values
within the original array and is not stable. Our version recurses on the
smaller partition and loops over the larger one, limiting recursion space
to $O(\log n)$. These timings illustrate the expected trends; they do not
prove the complexity bounds.

\end{document}
